\documentclass[11pt]{article}
\usepackage[margin=1in]{geometry}
\usepackage{amsmath,amssymb,amsthm}
\usepackage{physics}
\usepackage{booktabs}
\usepackage{hyperref}
\usepackage{cleveref}
\usepackage{algorithm}
\usepackage{algpseudocode}

\newtheorem{theorem}{Theorem}
\newtheorem{proposition}{Proposition}
\newtheorem{remark}{Remark}

\title{Numerical Demonstration: Solving the 1D Heat Equation\\
via Hybrid CV-DV LCHS}
\author{}
\date{}

\begin{document}
\maketitle

\begin{abstract}
We present a self-contained numerical demonstration of the hybrid
continuous-variable/discrete-variable (CV-DV) linear combination of Hamiltonian
simulation (LCHS) method applied to the one-dimensional heat equation.
We detail every mathematical step from PDE discretization through the LCHS
integral identity, state preparation via position-eigenbasis encoding, joint
Hamiltonian evolution, and post-selected measurement extraction.
The implementation achieves $0.2\%$ relative error at Fock dimension $N=300$
against the classical matrix exponential solution.
\end{abstract}

%% ====================================================================
\section{Problem Setup: Discretized 1D Heat Equation}
%% ====================================================================

Consider the one-dimensional heat equation on a rod of length $L$ with
homogeneous Dirichlet boundary conditions:
\begin{equation}\label{eq:heat-pde}
  \frac{\partial u}{\partial t} = \alpha \frac{\partial^2 u}{\partial x^2},
  \qquad u(0,t)=u(L,t)=0, \qquad u(x,0)=f(x),
\end{equation}
where $\alpha > 0$ is the thermal diffusivity.

\subsection{Spatial Discretization}

We discretize the spatial domain into $N_x + 1 = 5$ equally spaced points
with grid spacing $h = L/(N_x+1)$.  The boundary points $x_0$ and $x_5$ are
fixed at zero, leaving four interior unknowns
$\vb{u}(t) = (u_1(t),\,u_2(t),\,u_3(t),\,u_4(t))^\top$.
The second-order central difference approximation gives
\begin{equation}
  \frac{\partial^2 u}{\partial x^2}\bigg|_{x_j}
  \approx \frac{u_{j+1} - 2u_j + u_{j-1}}{h^2},
\end{equation}
leading to the semidiscrete ODE system
\begin{equation}\label{eq:ode}
  \dv{\vb{u}}{t} = -A\,\vb{u}(t), \qquad \vb{u}(0) = \vb{u}_0,
\end{equation}
where
\begin{equation}\label{eq:A-matrix}
  A = \frac{\alpha}{h^2}\,T, \qquad
  T = \begin{bmatrix}
    2 & -1 & 0 & 0 \\
    -1 & 2 & -1 & 0 \\
    0 & -1 & 2 & -1 \\
    0 & 0 & -1 & 2
  \end{bmatrix}.
\end{equation}
The exact solution is $\vb{u}(t) = e^{-tA}\,\vb{u}_0$.

\subsection{Pauli Decomposition}

Since $\vb{u} \in \mathbb{R}^4 \cong (\mathbb{C}^2)^{\otimes 2}$, we encode
the solution on two qubits.  The tridiagonal matrix admits the Pauli decomposition
\begin{equation}\label{eq:pauli}
  T = 2\,I_0 I_1 - I_0 X_1 - \tfrac{1}{2}(X_0 X_1 + Y_0 Y_1),
\end{equation}
so that $A = \frac{\alpha}{h^2}\bigl(2\,II - IX - \frac{1}{2}XX - \frac{1}{2}YY\bigr)$.

\subsection{Cartesian Decomposition}

Since $A$ is real symmetric, its Cartesian decomposition is trivial:
\begin{equation}\label{eq:cartesian}
  A = L + iH, \qquad L = \frac{A+A^\dagger}{2} = A, \qquad
  H = \frac{A-A^\dagger}{2i} = 0.
\end{equation}
The vanishing of $H$ means the entire dynamics is dissipative with no
Hamiltonian (unitary) part, and the LCHS circuit reduces to a single
CV-DV interaction term with no Trotter splitting needed.

\subsection{Numerical Parameters}

For the demonstration, we set:
\begin{center}
\begin{tabular}{lll}
  \toprule
  Parameter & Symbol & Value \\
  \midrule
  Thermal diffusivity & $\alpha$ & $1.0$ \\
  Grid spacing & $h$ & $1.0$ \\
  Total time & $T$ & $1.0$ \\
  Time steps & $n_{\mathrm{steps}}$ & $50$ \\
  Kernel parameter & $\beta$ & $0.8$ \\
  Fock dimension & $N$ & $100$--$300$ \\
  Initial state & $\ket{u_0}$ & $\ket{01}$ \\
  \bottomrule
\end{tabular}
\end{center}
The initial condition $\ket{u_0}=\ket{01}$ corresponds to
$\vb{u}_0 = (0,\,1,\,0,\,0)^\top$, i.e., a unit pulse at the second grid point.

%% ====================================================================
\section{The LCHS Integral Identity}
%% ====================================================================

The LCHS framework expresses the matrix exponential as a continuous integral
over unitary operators.

\begin{theorem}[CV-DV LCHS {\cite[Thm.~1]{das2025}}]\label{thm:lchs}
Let $A = L + iH$ with $L \succeq 0$.  Then
\begin{equation}\label{eq:lchs-integral}
  e^{-tA} = \int_{\mathbb{R}} g(k)\, e^{-it(kL + H)}\,\dd{k},
\end{equation}
where the kernel function is
\begin{equation}\label{eq:kernel}
  g(k) = \frac{f(k)}{1-ik}, \qquad
  f(k) = \frac{e^{2^\beta}}{2\pi\, e^{(1+ik)^\beta}},
\end{equation}
for a user-defined parameter $\beta \in (0,1)$.
\end{theorem}

Equivalently, combining constants:
\begin{equation}\label{eq:g-combined}
  g(k) = \frac{1}{C_\beta\, e^{(1+ik)^\beta}\,(1-ik)},
  \qquad C_\beta = 2\pi\, e^{-2^\beta}.
\end{equation}

\subsection{Verification of the Integral}

For our heat equation with $H=0$, \cref{eq:lchs-integral} becomes
\begin{equation}\label{eq:lchs-heat}
  e^{-tA} = \int_{-\infty}^{\infty} g(k)\, e^{-itkA}\,\dd{k}.
\end{equation}
Since $A$ is diagonalizable as $A = Q\Lambda Q^\top$ with
$\Lambda = \operatorname{diag}(\lambda_1,\ldots,\lambda_4)$, this reduces to
applying the scalar identity
\begin{equation}\label{eq:scalar-lchs}
  e^{-t\lambda} = \int_{-\infty}^{\infty} g(k)\, e^{-itk\lambda}\,\dd{k}
\end{equation}
independently to each eigenvalue $\lambda_j \geq 0$.
Numerical quadrature with sufficiently many points confirms
\cref{eq:scalar-lchs} to machine precision.

%% ====================================================================
\section{Quantum Circuit: CV-DV Implementation}
%% ====================================================================

The LCHS circuit uses a single bosonic (CV) mode as an ancilla coupled to
two DV qubits.  The protocol has three stages.

\subsection{Joint State and Evolution}

Prepare the joint initial state
\begin{equation}\label{eq:joint-init}
  \ket{\Psi_0} = \ket{\psi}_{\mathrm{osc}} \otimes \ket{u_0}_q,
\end{equation}
where $\ket{\psi}_{\mathrm{osc}}$ is a CV state encoding the kernel (see
\cref{sec:state-prep}), and $\ket{u_0}_q = \ket{01}$ is the initial
qubit state.

The joint Hamiltonian is
\begin{equation}\label{eq:joint-ham}
  \mathcal{H} = \hat{x} \otimes A,
\end{equation}
where $\hat{x} = (\hat{a}+\hat{a}^\dagger)/2$ is the position operator of the
oscillator.  The time evolution operator for a step $\delta t$ is
\begin{equation}\label{eq:U-step}
  U(\delta t) = e^{-i\,\delta t\,\hat{x}\otimes A}.
\end{equation}
Since $H=0$, this is the \emph{only} term---no Trotter decomposition is
required.

\subsection{Post-Selection and Solution Extraction}

After evolving to time $t$, we extract the qubit state by projecting the
oscillator onto a post-selection state $\bra{\phi}_{\mathrm{osc}}$:
\begin{equation}\label{eq:postselect}
  \vb{v}(t) = \bigl(\bra{\phi}_{\mathrm{osc}} \otimes I_q\bigr)\,
  U(t)\,\ket{\Psi_0}.
\end{equation}
By \cref{thm:lchs}, the key identity connecting the quantum and classical
solutions is
\begin{equation}\label{eq:sandwich}
  \bra{\phi}\, e^{-it\hat{x}\otimes A}\, \ket{\psi}
  = \biggl[\int_{\mathbb{R}} \phi^*(x)\,\psi(x)\, e^{-itxA}\,\dd{x}\biggr].
\end{equation}
If the CV states are engineered so that
$\phi^*(x)\,\psi(x) = g(x)$ (in the position representation), then
\cref{eq:sandwich} reproduces \cref{eq:lchs-heat} exactly, yielding
\begin{equation}
  \vb{v}(t) = e^{-tA}\,\vb{u}_0.
\end{equation}
In practice, the extracted vector $\vb{v}(t)$ carries a normalization factor
$\mathcal{N}$ from the state preparation, so the physical solution is recovered as
\begin{equation}\label{eq:normalize}
  \vb{u}_{\mathrm{LCHS}}(t) = \frac{\vb{v}(t)}{\mathcal{N}},
  \qquad \mathcal{N} = \norm{\vb{v}(0)}.
\end{equation}

%% ====================================================================
\section{State Preparation}\label{sec:state-prep}
%% ====================================================================

The central challenge is to prepare oscillator states
$\ket{\psi}$ and $\ket{\phi}$ whose position-space overlap reproduces the
LCHS kernel: $\phi^*(x)\,\psi(x) = g(x)$.

\subsection{Paper Method: Squeezed-Fock Expansion}

The original approach~\cite{das2025} chooses $\ket{\phi}$ as a squeezed vacuum
with squeezing parameter $r$:
\begin{equation}
  \phi_r(x) = \Bigl(\frac{2}{\sigma^2\pi}\Bigr)^{1/4} e^{-x^2/\sigma^2},
  \qquad \sigma = e^r.
\end{equation}
The corresponding preparation state must then satisfy
\begin{equation}\label{eq:psi-r}
  \psi_r(x) = \mathcal{N}\, g(x)\, e^{x^2/\sigma^2},
\end{equation}
which is expanded in the squeezed Fock basis
$\ket{\psi_r} = S(r') \sum_{n} C_n \ket{n}$ with coefficients
\begin{equation}\label{eq:Cn}
  C_n = \sqrt{\frac{\sigma}{\sigma'}}\,\frac{1}{\sqrt{2^n n!}}
  \int_{-\infty}^{\infty} H_n\!\bigl(x/\sigma'\bigr)\, g(x)\,
  e^{-\gamma x^2}\,\dd{x},
\end{equation}
where $\gamma = e^{-2r'} - e^{-2r} = \sigma'^{-2} - \sigma^{-2}$,
$\sigma' = e^{r'}$, and $r' < r$.

\begin{remark}[Practical limitation]
The factor $e^{x^2/\sigma^2}$ in \cref{eq:psi-r} makes
$\psi_r(x) \notin L^2(\mathbb{R})$ unless $\sigma$ is large enough that
$e^{x^2/\sigma^2} \approx 1$ over the support of $g(x)$.
Since $|g(x)|$ has support out to $|x| \sim 10$, one needs
$\sigma \gg 10$, i.e., $r \gg 2.3$.  But then the mean Fock occupation
$\langle \hat{n} \rangle = \sinh^2(r) \gg N$ for any practical Fock
cutoff $N$.  This makes the squeezed-Fock approach impractical for
PDE applications where $A$ enters through $\hat{x} \otimes A$.
\end{remark}

\subsection{Practical Method: Position-Eigenbasis Encoding}

We introduce an alternative approach that avoids squeezed states entirely
by working directly in the eigenbasis of $\hat{x}$.

\subsubsection{Truncated Position Operator}

In the $N$-dimensional Fock space, the position operator is
\begin{equation}\label{eq:x-fock}
  \hat{x} = \frac{\hat{a} + \hat{a}^\dagger}{2},
  \qquad (\hat{x})_{mn} = \frac{1}{2}\bigl(\sqrt{m+1}\,\delta_{m+1,n}
  + \sqrt{m}\,\delta_{m-1,n}\bigr).
\end{equation}
Diagonalizing this $N \times N$ matrix yields eigenvalues $\{x_j\}_{j=0}^{N-1}$
and eigenvectors $\{\ket{x_j}\}$:
\begin{equation}
  \hat{x}\,\ket{x_j} = x_j\,\ket{x_j}, \qquad
  \ket{x_j} = \sum_{n=0}^{N-1} V_{nj}\,\ket{n},
\end{equation}
where $V$ is the orthogonal matrix of eigenvectors.
The eigenvalues $x_j$ provide a natural quadrature grid on the
position axis, with spacing that becomes dense near $x=0$ and extends
to $|x| \sim \sqrt{2N}$.

\subsubsection{Integration Weights}

We assign midpoint-rule quadrature weights to each grid point:
\begin{equation}\label{eq:weights}
  w_j = \begin{cases}
    x_1 - x_0, & j=0, \\[4pt]
    \dfrac{x_{j+1}-x_{j-1}}{2}, & 1 \leq j \leq N-2, \\[4pt]
    x_{N-1} - x_{N-2}, & j = N-1.
  \end{cases}
\end{equation}

\subsubsection{State Construction}

We need $\phi^*(x_j)\,\psi(x_j) = w_j\,g(x_j)$ at each quadrature point so
that
\begin{equation}
  \sum_{j=0}^{N-1} \phi^*(x_j)\,\psi(x_j)\,e^{-itx_j A}
  \approx \int_{\mathbb{R}} g(x)\,e^{-itxA}\,\dd{x} = e^{-tA}.
\end{equation}
A natural factorization is:
\begin{equation}\label{eq:factorize}
  \boxed{%
    \psi(x_j) = \sqrt{w_j}\,g(x_j), \qquad
    \phi(x_j) = \sqrt{w_j}.}
\end{equation}
Converting from the position eigenbasis to the Fock basis:
\begin{equation}\label{eq:fock-conversion}
  \ket{\psi}_{\mathrm{osc}} = \mathcal{N}_\psi^{-1}
  \sum_{n=0}^{N-1}\biggl(\sum_{j=0}^{N-1} V_{nj}\,\sqrt{w_j}\,g(x_j)\biggr)
  \ket{n},
  \qquad
  \ket{\phi}_{\mathrm{osc}} =
  \sum_{n=0}^{N-1}\biggl(\sum_{j=0}^{N-1} V_{nj}\,\sqrt{w_j}\biggr)\ket{n},
\end{equation}
where $\mathcal{N}_\psi$ normalizes $\ket{\psi}$ to unit norm. The
post-selection state $\ket{\phi}$ is left unnormalized as its norm contributes
to the overall LCHS normalization factor $\mathcal{N}$ in \cref{eq:normalize}.

In matrix notation, writing $\vb{g} = (\sqrt{w_0}\,g(x_0),\ldots,
\sqrt{w_{N-1}}\,g(x_{N-1}))^\top$ and
$\vb{w} = (\sqrt{w_0},\ldots,\sqrt{w_{N-1}})^\top$:
\begin{equation}
  \psi_{\mathrm{Fock}} = V\,\vb{g}, \qquad
  \phi_{\mathrm{Fock}} = V\,\vb{w}.
\end{equation}

\subsubsection{Why This Works}

The eigenstates $\ket{x_j}$ satisfy the resolution of identity within the
truncated Fock space: $\sum_j \ket{x_j}\!\bra{x_j} = I_N$.  Therefore,
\begin{align}
  \bra{\phi}\,e^{-it\hat{x}\otimes A}\,\ket{\psi}
  &= \sum_j \bra{\phi}\ket{x_j}\bra{x_j}\,
     e^{-it\hat{x}\otimes A}\,\ket{\psi} \notag\\
  &= \sum_j \phi^*(x_j)\,e^{-itx_j A}\,
     \bigl(\textstyle\sum_k \braket{x_j}{x_k}\,\psi(x_k)\bigr) \notag\\
  &= \sum_j \phi^*(x_j)\,\psi(x_j)\,e^{-itx_j A} \notag\\
  &= \sum_j w_j\,g(x_j)\,e^{-itx_j A} \label{eq:quadrature-sum}\\
  &\approx \int_{\mathbb{R}} g(x)\,e^{-itxA}\,\dd{x}
   = e^{-tA}. \notag
\end{align}
The approximation error comes solely from discretizing the continuous
integral by the $N$-point quadrature rule defined by the eigenvalues of
$\hat{x}$.  As $N$ increases, the eigenvalues become denser and span a
wider range, so the quadrature converges.

%% ====================================================================
\section{Simulation Algorithm}
%% ====================================================================

\begin{algorithm}[h]
\caption{CV-DV LCHS for the 1D Heat Equation}
\label{alg:lchs}
\begin{algorithmic}[1]
\Require Kernel parameter $\beta$, Fock dimension $N$, time step $\delta t$,
         number of steps $n_{\mathrm{steps}}$, initial qubit state $\ket{u_0}$
\Ensure Approximate solution $\vb{u}_{\mathrm{LCHS}}(t)$ at each time step
\State Diagonalize $\hat{x}$ in the $N$-dimensional Fock space:
       $\hat{x} = V\,\mathrm{diag}(x_0,\ldots,x_{N-1})\,V^\top$
\State Compute weights $w_j$ via \cref{eq:weights}
\State Compute kernel values $g(x_j)$ via \cref{eq:g-combined}
\State Construct Fock-basis states via \cref{eq:fock-conversion}
\State Build joint Hamiltonian $\mathcal{H} = \hat{x}\otimes A$ in
       $\mathbb{C}^N \otimes \mathbb{C}^4$
\State Compute $U_{\mathrm{step}} = e^{-i\,\delta t\,\mathcal{H}}$
\State Initialize $\ket{\Psi} \gets \ket{\psi}_{\mathrm{osc}}\otimes\ket{u_0}_q$
\State Set normalization: $\mathcal{N} \gets
       \norm{(\bra{\phi}_{\mathrm{osc}}\otimes I_q)\ket{\Psi}}$
\For{$s = 0,\,1,\,\ldots,\,n_{\mathrm{steps}}$}
  \State Project: $\vb{v}(s\,\delta t) \gets
         \mathrm{Re}\bigl[(\bra{\phi}_{\mathrm{osc}}\otimes I_q)\ket{\Psi}\bigr]$
  \State Record: $\vb{u}_{\mathrm{LCHS}}(s\,\delta t) \gets \vb{v}(s\,\delta t)/\mathcal{N}$
  \State Evolve: $\ket{\Psi} \gets U_{\mathrm{step}}\,\ket{\Psi}$
\EndFor
\end{algorithmic}
\end{algorithm}

\begin{remark}[No Trotter error]
Since $H = 0$ for the heat equation, the joint evolution
$U(t) = e^{-it\hat{x}\otimes A}$ involves only a single exponential.
There is no need for Trotter--Suzuki decomposition and hence
\emph{zero Trotter error}.  The only source of error is the finite-dimensional
approximation to the continuous LCHS integral.
\end{remark}

%% ====================================================================
\section{Numerical Results}
%% ====================================================================

\subsection{Classical Reference}

The exact solution at $t=1$ with $\vb{u}_0 = (0,1,0,0)^\top$ and
$A = T$ (since $\alpha = h = 1$) is computed via matrix exponentiation:
\begin{equation}
  \vb{u}_{\mathrm{exact}}(1) = e^{-A}\,\vb{u}_0
  = \begin{pmatrix} 0.1864 \\ 0.3014 \\ 0.2126 \\ 0.0862 \end{pmatrix}.
\end{equation}
The eigenvalues of $A$ are
$\lambda_k = 2 - 2\cos\!\bigl(\frac{k\pi}{5}\bigr)$ for $k=1,2,3,4$,
giving $\lambda \approx 0.382,\; 1.382,\; 2.618,\; 3.618$.

\subsection{Convergence with Fock Dimension}

The position-eigenbasis method was tested at $\beta = 0.8$ for increasing
Fock dimension $N$.  The relative PDE error is defined as
\begin{equation}
  \varepsilon = \frac{\norm{\vb{u}_{\mathrm{LCHS}}(T) - \vb{u}_{\mathrm{exact}}(T)}}
                     {\norm{\vb{u}_{\mathrm{exact}}(T)}}.
\end{equation}

\begin{center}
\begin{tabular}{rcc}
  \toprule
  $N$ & Normalization $\mathcal{N}$ & Relative error $\varepsilon$ \\
  \midrule
  50  & $1.455$ & $1.17 \times 10^{-1}$ \\
  100 & $1.741$ & $4.80 \times 10^{-2}$ \\
  200 & $1.909$ & $4.69 \times 10^{-3}$ \\
  300 & $1.982$ & $2.06 \times 10^{-3}$ \\
  \bottomrule
\end{tabular}
\end{center}

The error decreases steadily as $N$ increases because:
\begin{enumerate}
  \item The eigenvalue range $[x_0, x_{N-1}]$ of the truncated $\hat{x}$ grows
        as $\sim \sqrt{2N}$, covering more of the kernel's support.
  \item The eigenvalue density increases, improving the quadrature accuracy
        in \cref{eq:quadrature-sum}.
  \item The normalization factor $\mathcal{N}$ approaches its continuum value,
        indicating that the overlap $\braket{\phi}{\psi}$ converges.
\end{enumerate}

\subsection{Component-wise Comparison at $N=300$}

\begin{center}
\begin{tabular}{ccc}
  \toprule
  Component & $u_{\mathrm{LCHS}}(T=1)$ & $u_{\mathrm{exact}}(T=1)$ \\
  \midrule
  $u_1$ & $0.18585$ & $0.18645$ \\
  $u_2$ & $0.30123$ & $0.30143$ \\
  $u_3$ & $0.21310$ & $0.21261$ \\
  $u_4$ & $0.08650$ & $0.08616$ \\
  \bottomrule
\end{tabular}
\end{center}

All four components agree to three significant figures.

\subsection{Error vs.\ Time}

The relative error is not monotonic in time.  At $t=0$, the error is exactly
zero by construction (the normalization $\mathcal{N}$ is chosen to enforce
$\vb{u}_{\mathrm{LCHS}}(0) = \vb{u}_0$).  The error peaks near $t \approx T/2$
and then decreases toward the final time, consistent with the quadrature
approximation being most challenging for intermediate $t$ values where
the oscillatory integrand $g(x)\,e^{-itxA}$ has the highest frequency content.

Representative error values at $N = 300$, $\beta=0.8$:
\begin{center}
\begin{tabular}{cc}
  \toprule
  Time $t$ & Relative error $\varepsilon(t)$ \\
  \midrule
  $0.000$ & $0$ \\
  $0.240$ & $4.20 \times 10^{-3}$ \\
  $0.500$ & $1.07 \times 10^{-2}$ \\
  $0.740$ & $8.23 \times 10^{-3}$ \\
  $1.000$ & $2.06 \times 10^{-3}$ \\
  \bottomrule
\end{tabular}
\end{center}

%% ====================================================================
\section{Analysis of the Kernel}
%% ====================================================================

The LCHS kernel $g(k)$ defined in \cref{eq:g-combined} has the following
properties at $\beta = 0.8$:

\begin{itemize}
  \item \textbf{Decay:} $|g(k)| \sim \exp\bigl(-\mathrm{Re}[(1+ik)^\beta]\bigr)
        \cdot |1-ik|^{-1}$ decays super-algebraically for large $|k|$.
        The effective support is approximately $|k| \lesssim 15$.
  \item \textbf{Peak:} $|g(k)|$ is peaked near $k=0$ with
        $g(0) = \frac{1}{C_\beta\,e} = \frac{e^{2^\beta}}{2\pi\,e}
        \approx 0.0640$.
  \item \textbf{Analyticity:} $g(k)$ is analytic in the strip
        $|\mathrm{Im}(k)| < 1$, which governs the exponential convergence
        rate of the Hermite expansion.
  \item \textbf{Normalization:}
        $\int_{\mathbb{R}} g(k)\,\dd{k} = 1$, since $e^{-0\cdot A} = I$.
\end{itemize}

For the position-eigenbasis method, what matters is that the eigenvalue
range $[x_0,\, x_{N-1}]$ of the truncated $\hat{x}$ must cover the essential
support of $g$.  At $N = 300$, we have $x_{N-1} \approx \sqrt{2 \cdot 300}
\approx 24.5$, which comfortably exceeds the kernel's support of
$|x| \lesssim 15$.

%% ====================================================================
\section{Comparison of State-Preparation Methods}
%% ====================================================================

\subsection{Squeezed-Fock (Paper Method)}

The paper's method requires $\psi_r(x) = \mathcal{N}\,g(x)\,e^{x^2/\sigma^2}$
to be square-integrable.  This demands $\sigma \gg$ kernel support, forcing
$r \gg 2.3$.  However, the mean Fock number of a squeezed vacuum is
$\langle\hat{n}\rangle = \sinh^2(r)$, which grows exponentially with $r$:

\begin{center}
\begin{tabular}{rrc}
  \toprule
  $r$ & $\langle\hat{n}\rangle$ & Feasible at $N=300$? \\
  \midrule
  $1.0$ & $1.4$ & Yes, but $\sigma = 2.7 \ll 10$ \\
  $2.0$ & $13.1$ & Yes, but $\sigma = 7.4 < 10$ \\
  $3.0$ & $100.4$ & Marginal ($\sigma = 20.1$) \\
  $5.0$ & $5{,}500$ & No ($N \gg 300$ needed) \\
  $8.0$ & $4.5\times 10^6$ & Completely impractical \\
  \bottomrule
\end{tabular}
\end{center}

The expansion coefficients $\{C_n\}$ are individually verified against
the paper's analytical results (Tables~1 and~2 therein) to 10-decimal-place
agreement.  However, the coefficient norms grow explosively with $N$:
$\norm{\vb{C}} \approx 2.8$ at $N=32$, $152$ at $N=64$, and $13{,}546$ at
$N=96$---indicating a divergent expansion that cannot be tamed at practical
Fock dimensions.

\subsection{Position-Eigenbasis (New Method)}

The position-eigenbasis method completely avoids the problematic
$e^{x^2/\sigma^2}$ factor:
\begin{itemize}
  \item No squeezing parameters $r,\,r'$ to tune.
  \item The only parameter is the Fock dimension $N$ (and the kernel
        parameter $\beta$).
  \item The error decreases monotonically with $N$.
  \item The method is numerically stable---no exponentially growing
        coefficients.
\end{itemize}

The tradeoff is that the position-eigenbasis encoding does not correspond
to a simple physical state-preparation circuit.  It serves as an
\emph{ideal reference} for the LCHS integral, isolating the quadrature
approximation error from all circuit-level errors.

%% ====================================================================
\section{Error Budget}
%% ====================================================================

The total simulation error can be decomposed as:
\begin{equation}\label{eq:error-budget}
  \varepsilon_{\mathrm{total}} = \varepsilon_{\mathrm{quad}}
  + \varepsilon_{\mathrm{Trotter}} + \varepsilon_{\mathrm{state}},
\end{equation}
where:
\begin{itemize}
  \item $\varepsilon_{\mathrm{quad}}$: Quadrature error from approximating
        $\int g(x)\,e^{-itxA}\,\dd{x}$ by a finite sum.  This is the
        \emph{only} error present in the position-eigenbasis method and
        decreases with $N$.
  \item $\varepsilon_{\mathrm{Trotter}}$: Trotter--Suzuki decomposition
        error.  For the heat equation, $H=0$, so
        $\varepsilon_{\mathrm{Trotter}} = 0$ exactly.
  \item $\varepsilon_{\mathrm{state}}$: State-preparation error from
        approximating $g(x)$ with a finite squeezed-Fock expansion.
        For the position-eigenbasis method,
        $\varepsilon_{\mathrm{state}} = 0$ by construction.
\end{itemize}

Thus the observed $0.2\%$ error at $N=300$ is purely
$\varepsilon_{\mathrm{quad}}$, demonstrating that the LCHS integral identity
itself is highly accurate once the continuous integral is well-approximated.

%% ====================================================================
\section{Resource Analysis: Classical vs.\ Quantum}
%% ====================================================================

A critical question for any quantum algorithm is: where does the quantum
advantage reside, and what classical computation is still required?
We analyze this carefully for CV-DV LCHS, distinguishing three phases
of the algorithm.

\subsection{Phase 1: Classical Preprocessing}

The CV states $\ket{\psi}$ and $\ket{\phi}$ encode the LCHS kernel
$g(x)$, whose functional form depends on $\beta$ alone.
However, the required Fock cutoff $N$ is \emph{not} purely a property
of the kernel---it also depends on the oscillatory resolution demands
of the integrand $g(x)\,e^{-itxA}$.  Specifically, the quadrature grid
(eigenvalues of $\hat{x}$) must be dense enough to resolve oscillations
at frequency $\sim t\|A\|$.  Thus:
\begin{equation}\label{eq:N-dependence}
  N = N(\beta,\; t\|A\|,\; \epsilon_{\mathrm{quad}}).
\end{equation}
$N$ depends on the product $t\|A\|$ (which sets the oscillation frequency)
and the target accuracy, but \emph{not on the system dimension $2^m$
directly}.  A larger PDE system (more grid points) increases $m$ but
does not change $\|A\|$ for fixed grid spacing $h$ and diffusivity $\alpha$.

For the position-eigenbasis method, the classical preprocessing is:
\begin{center}
\begin{tabular}{lll}
  \toprule
  Step & Operation & Cost \\
  \midrule
  Diagonalize $\hat{x}$ & Eigendecompose $N\times N$ tridiagonal matrix
  & $O(N^2)^*$ \\
  Compute kernel & Evaluate $g(x_j)$ at $N$ points & $O(N)$ \\
  Basis transform & Matrix--vector products $V\vb{g}$, $V\vb{w}$ & $O(N^2)$ \\
  \midrule
  \textbf{Total} & & $O(N^2)$ \\
  \bottomrule
\end{tabular}
\end{center}
{\small ${}^*$The truncated $\hat{x}$ is tridiagonal, so a specialized
eigensolver (e.g., divide-and-conquer) achieves $O(N^2)$.
A general dense eigensolver costs $O(N^3)$; our implementation uses
\texttt{numpy.linalg.eigh} which falls in this category.}

\medskip\noindent
For our demonstration, $N = 100$--$300$ suffices for sub-percent accuracy
at $t\|A\| \approx 3.6$.

\begin{remark}[Position-eigenbasis as a numerical reference]
The position-eigenbasis method is a \emph{discrete quadrature surrogate},
not a direct physical state-preparation protocol.  It computes idealized
Fock-basis amplitudes $\{\psi_n\}$ and $\{\phi_n\}$ classically, then
assumes these can be loaded into the oscillator.  This shifts the
difficulty from physical Gaussian operations to classical dense linear
algebra and subsequent quantum state loading.
The method therefore serves as an ideal numerical reference that
isolates the quadrature approximation error from all circuit-level
state-preparation errors.
\end{remark}

\subsection{Phase 2: Quantum State Loading}

Loading an arbitrary $N$-component Fock-basis state
$\ket{\psi} = \sum_{n=0}^{N-1}\psi_n\ket{n}$ into a physical oscillator
is itself a nontrivial task.  The cost depends on the hardware platform:

\begin{center}
\begin{tabular}{lll}
  \toprule
  Platform & Method & Gate count \\
  \midrule
  Superconducting cavity & SNAP + displacement~\cite{heeres2015} & $O(N)$ \\
  Trapped ion & Sequential sideband pulses & $O(N)$ \\
  Photonic & Programmable interferometer & $O(N\log N)$ \\
  \bottomrule
\end{tabular}
\end{center}
In all cases, the loading cost scales with the \emph{Fock cutoff $N$},
not with the system size $2^m$.
This is a one-time cost per run; the loaded states are consumed by the
subsequent quantum evolution.

For the paper's squeezed-Fock method, loading is simpler in principle
(apply $S(r')$ to a Fock superposition), but requires impractically large
squeezing as discussed in \cref{sec:state-prep}.

\subsection{Phase 3: Quantum Hamiltonian Evolution}

This is where the quantum advantage resides.
The joint evolution $U(t)=e^{-it\hat{x}\otimes A}$ acts on the composite
Hilbert space $\mathcal{H}_{\mathrm{osc}} \otimes \mathcal{H}_{\mathrm{DV}}$
of dimension $N \times 2^m$.  For a general $A$ with Cartesian decomposition
$A = L + iH$, the Trotter--Suzuki decomposition gives:

\paragraph{Gate complexity per Trotter step.}
Let $L = \sum_{k=1}^{N_L} \alpha_k P_k$ and $H = \sum_{k=1}^{N_H}\beta_k P_k'$
be the Pauli decompositions.  Each term $e^{-i\delta t\,\alpha_k(\hat{x}\otimes P_k)}$
is implemented by a conditional displacement (CD) gate sandwiched between
single-qubit basis rotations.  One Trotter step requires:
\begin{equation}
  N_{\mathrm{CD}} = N_L \text{ (hybrid CV-DV gates)}, \qquad
  N_{\mathrm{CNOT}} = 2\sum_{k=1}^{N_L+N_H}(w(P_k)-1)
  \text{ (DV-only gates)},
\end{equation}
where $w(P)$ is the Pauli weight (number of non-identity factors).

\paragraph{Scaling for the 1D heat equation.}
For the 1D heat equation on $d = 2^m$ grid points encoded on $m$ qubits,
the Pauli decomposition of the tridiagonal Laplacian contains
$N_L = O(d)$ terms (one per nearest-neighbor bond), each with Pauli weight
$w \leq 2$.  Therefore:
\begin{itemize}
  \item \textbf{Total gate count per Trotter step:} $O(d)$ CD gates and
        $O(d)$ CNOT gates.
  \item \textbf{Circuit depth per Trotter step:} $O(1)$.
        The even-odd bond splitting~\cite[Eq.~49--51]{das2025} groups all
        even bonds (mutually commuting) into one layer and all odd bonds
        into another, giving exactly 2 circuit layers regardless of $d$.
\end{itemize}
The distinction between gate count and circuit depth is critical:
the total number of gates scales linearly in $d$, but they execute in
$O(1)$ parallel layers.

\paragraph{Number of Trotter steps.}
The first-order Trotter error bound gives
\begin{equation}
  n = O\!\left(\frac{t^2}{\epsilon_t}\|\hat{x}_N\|\,
  \|\hat{L}\|_1\|\hat{H}\|_1\right),
\end{equation}
where $\|\hat{x}_N\| \sim \sqrt{N}$ on the truncated Fock space,
$\|\hat{L}\|_1 = \sum|\alpha_k|$, and $\|\hat{H}\|_1 = \sum|\beta_j|$.
For the heat equation ($H=0$), no Trotter steps are needed at all.

\paragraph{Total quantum resources.}
\begin{center}
\begin{tabular}{lll}
  \toprule
  Resource & Symbol & Scaling \\
  \midrule
  Oscillator modes & & $1$ (single ancilla) \\
  DV qubits & $m$ & $\lceil\log_2 d\rceil$ \\
  Fock cutoff & $N$ & $N(\beta,\,t\|A\|,\,\epsilon)$;
  see \cref{eq:N-dependence} \\
  Trotter steps & $n$ & $O\!\left(\frac{t^2\sqrt{N}\|\hat{L}\|_1
  \|\hat{H}\|_1}{\epsilon_t}\right)$ \\
  CV-DV gates per step & & $O(N_L)$ CD gates \\
  DV gates per step & & $O(\sum w(P_k))$ CNOT gates \\
  Circuit depth per step & & $O(1)$ for 1D nearest-neighbor \\
  \bottomrule
\end{tabular}
\end{center}

\subsection{Classical vs.\ Quantum Cost Comparison}

To solve $\vb{u}(t) = e^{-tA}\vb{u}_0$ for the 1D heat equation with
$d = 2^m$ grid points:
\begin{center}
\begin{tabular}{lll}
  \toprule
  & Classical & Quantum (CV-DV LCHS) \\
  \midrule
  State storage & $O(d)$ real numbers & $m$ qubits + 1 oscillator \\
  Per-step gate count & $O(d)$ (sparse mat-vec) & $O(d)$ (CD + CNOT) \\
  Per-step circuit depth & --- & $O(1)$ (parallelizable) \\
  Full solve & $O(d \cdot n_t)$ & $O(n \cdot d)$ gates \\
  Direct exponentiation & $O(d^3)$ or $O(d^2)$ & --- \\
  State prep (one-time) & --- & $O(N)$ oscillator gates \\
  Preprocessing & --- & $O(N^2)$ classical \\
  \bottomrule
\end{tabular}
\end{center}

\begin{remark}[Honest assessment for 1D]
For the \emph{1D} heat equation, the Laplacian is tridiagonal and
admits $O(d)$-cost sparse classical solvers.  The quantum per-step gate
count is also $O(d)$---there is \textbf{no asymptotic gate-count advantage}
in 1D with naive Pauli-term Trotterization.

The potential advantages of the quantum approach are:
\begin{enumerate}
  \item \textbf{Exponential state-space compression:}
        The classical solver stores a $d$-component vector; the quantum
        circuit encodes it in $m = \log_2 d$ qubits.
        Extracting full-state information requires $O(d)$ measurements,
        but problems requiring only expectation values (e.g., average
        temperature) can exploit this compression.
  \item \textbf{Constant circuit depth:}
        Each Trotter step has $O(1)$ depth for nearest-neighbor 1D
        Hamiltonians, regardless of $d$.  On parallel quantum hardware,
        this yields wall-clock time independent of system size.
  \item \textbf{Higher dimensions:}
        For $D$-dimensional PDEs on $d = n^D$ grid points, the Pauli
        decomposition structure and classical solver costs differ.
        Block-encoding or qubitization approaches can achieve
        $\tilde{O}(\mathrm{poly}(m))$ per-step cost, recovering
        exponential advantage over classical $O(d)$-per-step methods.
  \item \textbf{CV ancilla advantage over discrete LCU:}
        The single oscillator replaces $\lceil\log_2 M\rceil$ ancilla
        qubits that discrete LCHS would require, where
        $M \sim T\|L\|(\log 1/\epsilon)^{1+1/\beta}$ can be $\sim 10^6$.
\end{enumerate}
\end{remark}

\subsection{Resource Summary for the Toy Problem}

For our 4-point demonstration ($m=2$, $d=4$), all resources are trivially
small and there is no quantum advantage:
\begin{center}
\begin{tabular}{lll}
  \toprule
  Resource & Value & Notes \\
  \midrule
  \multicolumn{3}{l}{\textit{Classical preprocessing}} \\
  Diagonalize $\hat{x}$ & $300 \times 300$ eigenproblem & 0.01 sec \\
  Kernel evaluation & $g(x_j)$ at 300 points & negligible \\
  Basis transform & $300\times 300$ mat-vec & negligible \\
  \midrule
  \multicolumn{3}{l}{\textit{Quantum resources}} \\
  Oscillator modes & 1 & single CV ancilla \\
  DV qubits & 2 & $\lceil\log_2 4\rceil$ \\
  Fock cutoff $N$ & 100--300 & controls quadrature accuracy \\
  Pauli terms & 4 ($II$, $IX$, $XX$, $YY$) & from \cref{eq:pauli} \\
  Trotter steps & 0 & $H=0$ for heat equation \\
  Osc.\ state loading & $O(N)$ gates & SNAP+disp.\ or equivalent \\
  \midrule
  \multicolumn{3}{l}{\textit{Classical reference (no quantum)}} \\
  Matrix exponential & $4\times 4$ dense expm & negligible \\
  \bottomrule
\end{tabular}
\end{center}

The purpose of this toy problem is not to demonstrate quantum advantage
but to \emph{validate the mathematical framework}: confirming that the
LCHS integral identity, the CV-DV Hamiltonian evolution, and the
post-selection extraction all work correctly before scaling to larger
systems where classical simulation becomes intractable.

%% ====================================================================
\section{Discussion}
%% ====================================================================

\subsection{Summary of Findings}

This numerical demonstration establishes several key points:

\begin{enumerate}
  \item \textbf{The LCHS integral identity works.}  The continuous integral
        $\int g(k)\,e^{-itkA}\,\dd{k} = e^{-tA}$ is verified to high
        accuracy, confirming the theoretical foundation.

  \item \textbf{The bottleneck is state preparation.}  The paper's
        squeezed-Fock expansion requires impractically large squeezing
        parameters for PDE applications.  The position-eigenbasis encoding
        bypasses this, achieving $< 1\%$ error at moderate Fock dimensions.

  \item \textbf{No Trotter error for real symmetric $A$.}  The heat equation's
        Cartesian decomposition gives $H=0$, so the entire evolution is a
        single exponential $e^{-it\hat{x}\otimes A}$.

  \item \textbf{Convergence is controlled by $N$.}  The error decreases
        systematically as the Fock dimension increases, with the
        eigenvalue range of $\hat{x}$ growing as $\sqrt{2N}$ to cover
        the kernel support.
\end{enumerate}

\subsection{Where Is the Quantum Advantage?}

As analyzed in the resource comparison, the quantum advantage of
CV-DV LCHS is \emph{nuanced and problem-dependent}:

\begin{itemize}
  \item For the 1D heat equation with naive Pauli-term Trotterization,
        there is no asymptotic gate-count advantage over sparse classical
        solvers: both scale as $O(d)$ per step.
  \item The advantages lie in \textbf{state compression}
        ($m = \log_2 d$ qubits vs.\ $d$ classical variables),
        \textbf{constant circuit depth} per Trotter step (for
        nearest-neighbor Hamiltonians), and the \textbf{CV ancilla}
        replacing $O(\log M)$ discrete ancilla qubits.
  \item For higher-dimensional PDEs or with advanced simulation
        strategies (block-encoding, qubitization), the per-step cost can
        be reduced to $\tilde{O}(\mathrm{poly}(m))$, recovering exponential
        advantage.
\end{itemize}

The position-eigenbasis method developed here is a \emph{classical
numerical tool} for validating the LCHS protocol.  It proves that the
quadrature approximation error is small ($0.2\%$ at $N=300$) independently
of any circuit-level concerns.  Translating this to a physical
implementation requires either:
\begin{itemize}
  \item finding efficient Gaussian-operation circuits that approximate the
        idealized Fock amplitudes (an open problem), or
  \item using hardware-native arbitrary state preparation (e.g., SNAP gates
        on superconducting cavities), accepting the $O(N)$ gate overhead.
\end{itemize}
The state-loading cost $O(N)$ depends on the Fock cutoff
(which scales with $t\|A\|$ and $\epsilon$) but \emph{not on the system
dimension $2^m$}.

\begin{thebibliography}{9}
\bibitem{das2025}
E.~R.~Das, M.~Zheng, R.~Dutta, A.~Li, and Y.~Liu,
``Hybrid Oscillator-Qubit Linear Combination of Hamiltonian Simulation,''
2025.

\bibitem{an2023}
D.~An, J.-P.~Liu, and L.~Lin,
``Linear combination of Hamiltonian simulation for non-unitary dynamics
with optimal state preparation cost,''
\textit{Phys.\ Rev.\ Lett.}, vol.~131, 150603, 2023.

\bibitem{an2023theory}
D.~An and L.~Lin,
``Quantum linear system solver based on time-optimal adiabatic quantum
computing and quantum approximate optimization algorithm,''
\textit{ACM Trans.\ Quantum Comput.}, 2023.

\bibitem{heeres2015}
R.~W.~Heeres, B.~Vlastakis, E.~Holland, S.~Krastanov, V.~V.~Albert,
L.~Frunzio, L.~Jiang, and R.~J.~Schoelkopf,
``Cavity state manipulation using photon-number selective phase gates,''
\textit{Phys.\ Rev.\ Lett.}, vol.~115, 137002, 2015.
\end{thebibliography}

\end{document}
